% !TeX root = ../qm_chap3.tex

\begin{frame}{本章概览：算符理论的完整框架}
  \begin{iqbitemize}
    \item \textbf{算符基础}（Operator Basics）：定义、线性算符、常见算符实例
    \item \textbf{算符代数}（Operator Algebra）：加法、乘法、幂运算、算符函数
    \item \textbf{对易关系}（Commutation Relations）：对易子及其物理意义
    \item \textbf{本征值问题}（Eigenvalue Problems）：本征方程、可观测量、完备性
  \end{iqbitemize}

  \iqbbigsep

  \iqbgreenbox{核心思想}{算符是量子力学的语言，将物理量转化为数学对象}
\end{frame}
